\section{USAMO 2013}
        \begin{enumerate}
        	\item (\textit{USAMO 2013}) In triangle $ABC$, points $P,Q,R$ lie on sides $BC,CA,AB$ respectively.  Let $\omega_A$, $\omega_B$, $\omega_C$ denote the circumcircles of triangles $AQR$, $BRP$, $CPQ$, respectively.  Given the fact that segment $AP$ intersects $\omega_A$, $\omega_B$, $\omega_C$ again at $X,Y,Z$ respectively, prove that $YX/XZ=BP/PC$.


 



\textbf{Solution 1}


                    \begin{figure}[H]
                        \centering
                        \includegraphics[width=0.5\textwidth]{Solution/USAMO_I_2013_1_Solution_Figure_1.png}
                    \end{figure}
                    

 In this solution, all lengths and angles are directed.

 Firstly, it is easy to see by that $\omega_A, \omega_B, \omega_C$ concur at a point $M$ (the Miquel point). Let $XM$ meet $\omega_B, \omega_C$ again at $D$ and $E$, respectively. Then by Power of a Point, we have 
 \[XM \cdot XE = XZ \cdot XP \quad\text{and}\quad XM \cdot XD = XY \cdot XP\] Thusly 
 \[\frac{XY}{XZ} = \frac{XD}{XE}\] But we claim that $\triangle XDP \sim \triangle PBM$. Indeed, 
 \[\measuredangle XDP = \measuredangle MDP = \measuredangle MBP = - \measuredangle PBM\] and 
 \[\measuredangle DXP = \measuredangle MXY = \measuredangle MXA = \measuredangle MRA = 180^\circ - \measuredangle MRB = \measuredangle MPB = -\measuredangle BPM\]Therefore, $\frac{XD}{XP} = \frac{PB}{PM}$. Analogously we find that $\frac{XE}{XP} = \frac{PC}{PM}$ and we are done.

 
courtesy v\_enhance, minor clarification by integralarefun

 



\textbf{Solution 2}

\href{https://www.flickr.com/photos/127013945@N03/14800492500/lightbox/}{Diagram}Refer to the Diagram link.

 By Miquel's Theorem, there exists a point at which $\omega_A, \omega_B, \omega_C$ intersect. We denote this point by $M.$ Now, we angle chase:
 \[\angle YMX = 180^{\circ} - \angle YXM - \angle XYM\]
 \[= 180^{\circ} - \angle AXM - \angle PYM\]
 \[= \left(180^{\circ} - \angle ARM\right) - \angle PRM\]
 \[= \angle BRM - \angle PRM\]
 \[= \angle BRP = \angle BMP.\]In addition, we have
 \[\angle ZMX = 180^{\circ} - \angle MZY - \angle ZYM - \angle YMX\]
 \[= 180^{\circ} - \angle MZP - \angle PYM - \angle BMP\]
 \[= 180^{\circ} - \angle MCP - \angle PBM - \angle BMP\]
 \[= \left(180^{\circ} - \angle PBM - \angle BMP\right) - \angle MCP\]
 \[= \angle BPM - \angle MCP\]
 \[= 180^{\circ} - \angle MPC - \angle MCP\]
 \[= \angle CMP.\]Now, by the Ratio Lemma, we have
 \[\frac{XY}{XZ} = \frac{MY}{MZ} \cdot \frac{\sin \angle YMX}{\sin \angle ZMX}\]
 \[= \frac{\sin \angle YZM}{\sin \angle ZYM} \cdot \frac{\sin \angle BMP}{\sin \angle CMP}\] (by the Law of Sines in $\triangle MZY$)
 \[= \frac{\sin \angle PZM}{\sin \angle PYM} \cdot \frac{\sin \angle BMP}{\sin \angle CMP}\]
 \[= \frac{\sin \angle PCM}{\sin \angle PBM} \cdot \frac{\sin \angle BMP}{\sin \angle CMP}\]
 \[= \frac{MB}{MC} \cdot \frac{\sin \angle BMP}{\sin \angle CMP}\] (by the Law of Sines in $\triangle MBC$)
 \[= \frac{PB}{PC}\] by the Ratio Lemma.The proof is complete.

 



\textbf{Solution 3}

Use directed angles modulo $\pi$.

 Lemma. $\angle{XRY} \equiv \angle{XQZ}.$

 Proof. 
 \[\angle{XRY} \equiv \angle{XRA} - \angle{YRA} \equiv \angle{XQA} + \angle{YRB} \equiv \angle{XQA} + \angle{CPY} = \angle{XQA} + \angle{AQZ} = \angle{XQZ}.\]

 Now, it follows that (now not using directed angles)
 \[\frac{XY}{YZ} = \frac{\frac{XY}{\sin \angle{XRY}}}{\frac{YZ}{\sin \angle{XQZ}}} = \frac{\frac{RY}{\sin \angle{RXY}}}{\frac{QZ}{\sin \angle{QXZ}}} = \frac{BP}{PC}\]using the facts that $ARY$ and $APB$, $AQZ$ and $APC$ are similar triangles, and that $\frac{RA}{\sin \angle{RXA}} = \frac{QA}{\sin \angle{QXA}}$ equals twice the circumradius of the circumcircle of $AQR$.

 



\textbf{Solution 4}

We will use some construction arguments to solve the problem. Let $\angle BAC=\alpha,$ $\angle ABC=\beta,$ $\angle ACB=\gamma,$ and let $\angle APB=\theta.$ We construct lines through the points $Q,$ and $R$ that intersect with $\triangle ABC$ at the points $Q$ and $R,$ respectively, and that intersect each other at $T.$ We will construct these lines such that $\angle CQV=\angle ARV=\theta.$

 
Now we let the intersections of $AP$ with $RV$ and $QU$ be $Y'$ and $Z',$ respectively. This construction is as follows.
                    \begin{figure}[H]
                        \centering
                        \includegraphics[width=0.5\textwidth]{Solution/USAMO_I_2013_1_Solution_Figure_2.png}
                    \end{figure}
                    

 We know that $\angle BRY'=180^\circ-\angle ARY'=180^\circ-\theta.$ Hence, we have,
 \begin{align*} \angle BRY'+\angle BPY' &=180^\circ-\theta+\theta\\ &=180^\circ. \end{align*}


 Since the opposite angles of quadrilateral $RY'PB$ add up to $180^\circ,$ it must be cyclic. Similarly, we can also show that quadrilaterals $CQZ'P,$ and $AQTR$ are also cyclic.

 Since points $Y'$ and $Z'$ lie on $AP,$ we know that,
 \[Y'=\omega_B\cap AP\]and that
 \[Z'=\omega_C\cap AP.\]

 Hence, the points $Y'$ and $Z'$ coincide with the given points $Y$ and $Z,$ respectively.

 Since quadrilateral $AQTR$ is also cyclic, we have,
 \begin{align*} \angle Y'TZ' &=180^\circ-\angle RTQ\\ &=180^\circ-(180^\circ-\angle RAQ)\\ &=\angle RAQ\\ &=\alpha. \end{align*}


 Similarly, since quadrilaterals $CQZ'P,$ and $AQTR$ are also cyclic, we have,
 \begin{align*} \angle TY'Z' &=180^\circ-\angle RY'P\\ &=180^\circ-(180^\circ-\angle RBP)\\ &=\angle RBP\\ &=\beta, \end{align*}
and,
 \begin{align*} \angle Y'Z'T &=180^\circ-\angle PZ'Q\\ &=180^\circ-(180^\circ-\angle PCQ)\\ &=\angle PCQ\\ &=\gamma. \end{align*}


 Since these three angles are of $\triangle TY'Z',$ and they are equal to corresponding angles of $\triangle ABC,$ by AA similarity, we know that $\triangle TY'Z'\sim \triangle ABC.$

 We now consider the point $X=\omega_c\cap AC.$ We know that the points $A,$ $Q,$ $T,$ and $R$ are concyclic. Hence, the points $A,$ $T,$ $X,$ and $R$ must also be concyclic.

 Hence, quadrilateral $AQTX$ is cyclic.

 
                    \begin{figure}[H]
                        \centering
                        \includegraphics[width=0.5\textwidth]{Solution/USAMO_I_2013_1_Solution_Figure_3.png}
                    \end{figure}
                    

 Since the angles $\angle ART$ and $\angle AXT$ are inscribed in the same arc $\overarc{AT},$ we have,
 \begin{align*} \angle AXT &=\angle ART\\ &=\theta. \end{align*}


 Consider by this result, we can deduce that the homothety that maps $ABC$ to $TY'Z'$ will map $P$ to $X.$ Hence, we have that,
 \[Y'X/XZ'=BP/PC.\]

 Since $Y'=Y$ and $Z'=Z$ hence,
 \[YX/XZ=BP/PC,\]

 as required.

 



\textbf{Solution 5 (Simple Rotational Homotethy)}


                    \begin{figure}[H]
                        \centering
                        \includegraphics[width=0.5\textwidth]{Solution/USAMO_I_2013_1_Solution_Figure_4.png}
                    \end{figure}
                    

 We begin again by noting that the three circumcircles intersect at point $M$ by Miquel's theorem. In addition, we state that the angle $\angle MQC = \alpha$, hence $\angle MPC = \angle MZC = 180 - \alpha$, as well as $\angle AQM$, from which follows that $\angle ARM = \alpha$, so $\angle BRM = 180 - \alpha$, and $\angle BPM = \alpha$. We shall prove that the triangles $\triangle MZC$ and $\triangle XMP$ and $\triangle YMB$ are similar, which will imply a rotational homotethy with angle $\angle APB$ about the point $M$, that takes $Y,X,Z$ to $B,P,C$, thus proving the problem. (In essence, just imagine we rotate $YMZX$ around M and lengthen things out and get $MCPB$ - the ratios will remain identical.)

 We do this by angle chasing. Denote angle $\angle ZMP = \beta$. From the angles labeled before, we now know $\angle ZPB = \alpha-\beta = \angle ZMC$. In addition, $\angle ZPM = \angle ZCM$. So the angles in triangle $\triangle ZMC$  are $\alpha-\beta$, $\beta$ and thus $180-\alpha$, with $\alpha-\beta$ at the point $M$. In addition, $\angle YRM = \angle YPM = \beta$, so $\angle ARY = \alpha-\beta$, so $\angle YMB = \alpha - \beta$. Since $\angle YBM = \angle YPM = \beta$, the triangle $\triangle YMB$ has angles $\alpha-\beta, \beta, 180-\alpha$ also, with $\alpha-\beta$ at $M$. Finally, In triangle $\triangle XPM$, we already know the angle $\angle XPM$ to be $\beta$; we also can find that $\angle AQM = 180-\alpha$, so $\angle AXM = \alpha$, so $\angle MXP = 180-\alpha$, so $\angle XMP = \alpha-\beta$. Thus, the three triangles are similar have a common point $M$, which proves that there is a rotational homotethy around $M$ that maps $Y, X, Z$ to $B, P, C$ as desired. 

 \textit{Solution by SimilarTriangle.}

 


	\item (\textit{USAMO 2013}) For a positive integer $n\geq 3$ plot $n$ equally spaced points around a circle.  Label one of them $A$, and place a marker at $A$.  One may move the marker forward in a clockwise direction to either the next point or the point after that.  Hence there are a total of $2n$ distinct moves available; two from each point.  Let $a_n$ count the number of ways to advance around the circle exactly twice, beginning and ending at $A$, without repeating a move.  Prove that $a_{n-1}+a_n=2^n$ for all $n\geq 4$.


 



\textbf{Solution 1}

We label the points in clockwise order as $1,2,3,\ldots,n$, where point $A$ is the same as point $1$.  We start and end at point $1$, and we must cross over it, either by visiting it again, or else by making the move from point $n$ to point $2$. We interpret each of these cases in terms of tiling. In each move, we either move one or two points clockwise, so we can think of each move as a $1\times 1$ or $1\times 2$ tile. If the point $1$ is visited in the middle, then the first cycle around the circle can be thought of as a tiling of a $1\times n$ board, and the second cycle around the circle can also be thought of as a $1\times n$ board. We place this second board directly below the first board. Therefore, in this first case, we wish to find the number of tilings of two $1\times n$ boards, and to guarantee that no move is repeated, we cannot have two tiles of the same type lying directly atop each other in the $2\times n$ board. Suppose there are $u_n$ such tilings. It can easily be computed that $u_3=2$ and $u_4=6$ as shown below.
                    \begin{figure}[H]
                        \centering
                        \includegraphics[width=0.5\textwidth]{Solution/USAMO_I_2013_2_Solution_Figure_1.png}
                    \end{figure}
                    In the second case, where we pass by point $1$ by moving from point $n$ to point $2$, we can similarly think about it in terms of tiling two rows of $1\times n$ boards, but we remove the last square in the first row and the first square in the second row to make sure that we jump from point $n$ to point $2$. Suppose that we can tile such boards in $s_n$ ways. It can be easily computed that $s_3=3$ and $s_4=5$ as shown below.
                    \begin{figure}[H]
                        \centering
                        \includegraphics[width=0.5\textwidth]{Solution/USAMO_I_2013_2_Solution_Figure_2.png}
                    \end{figure}
                    Since these are the only two possible cases, we see that 
 \[a_n=u_n+s_n\tag{1}.\]

 For the sake of convenience in determining recurrence relations, we define another type of board with two $1\times n$ boards where a specific corner is removed (without loss of generality, we place this in the lower left hand corner). Let $t_n$ be the number of ways to tile such a board without placeing two of the same type of tile atop each other. Once again, we compute that $t_3=3$ and $t_4=5$ as shown below.
                    \begin{figure}[H]
                        \centering
                        \includegraphics[width=0.5\textwidth]{Solution/USAMO_I_2013_2_Solution_Figure_3.png}
                    \end{figure}
                    We can determine reccurence relations for $s_n$, $t_n$, and $u_n$ in terms of each other. For $s_n$, note that a tiling can end in one of the three following ways such that the rest of the board can be tiled without restriction (the placed tiles are shaded).
                    \begin{figure}[H]
                        \centering
                        \includegraphics[width=0.5\textwidth]{Solution/USAMO_I_2013_2_Solution_Figure_4.png}
                    \end{figure}
                    In the first, second, and third cases, we see that we can tile the rest of the board in $t_{n-1}$, $t_{n-2}$, and $s_{n-2}$ ways, respectively. Hence for $n\ge 5$, we see that
 \[s_n=t_{n-1}+t_{n-2}+s_{n-2}\tag{2}.\]For $t_n$, note that a tiling can end in one of the three following ways such that the rest of the board can be tiled without restriction.
                    \begin{figure}[H]
                        \centering
                        \includegraphics[width=0.5\textwidth]{Solution/USAMO_I_2013_2_Solution_Figure_5.png}
                    \end{figure}
                    In the first, second, and third cases, we see that we can tile the rest of the board in $s_{n-1}$, $s_{n-2}$, and $t_{n-2}$ ways, respectively. Hence for $n\ge 5$, we see that
 \[t_n=s_{n-1}+s_{n-2}+t_{n-2}\tag{3}.\]For $u_n$, note that a tiling can end in one of the two following ways such that the rest of the board can be tiled without restriction.
                    \begin{figure}[H]
                        \centering
                        \includegraphics[width=0.5\textwidth]{Solution/USAMO_I_2013_2_Solution_Figure_6.png}
                    \end{figure}
                    

 In either case, we are left with a board with a corner removed, hence we can tile the rest of the board in $t_{n-1}$  ways in each case. Hence for $n\ge 4$, we see that
 \[u_n=2t_{n-1}\tag{4}.\]Subtracting (3) from (2), we find that
 \[s_n-t_n=s_{n-1}-t_{n-1}.\]Therefore, if $s_{n-1}=t_{n-1}$, then $s_n=t_n$. Since $s_3=t_3$ and $s_4=t_4$, we see that $s_n=t_n$ for all $n\ge 3$. Therefore, (2) can be rewritten as
 \[s_n=s_{n-1}+2s_{n-2}\tag{5},\]and (4) can be rewritten as
 \[u_n=2s_{n-1}\tag{6}.\]Now by (1), we know that
 \[a_{n}+a_{n-1}=u_n+s_n+u_{n-1}+s_{n-1}.\tag{7}\]In particular, $a_4+a_3=6+5+2+3=2^4$, so the statement is true for $n=4$. Then by (7), and then substituting (5) and (6) (where these are valid for $n\ge 5$), we find
 \begin{align*} a_n+a_{n-1}&=u_n+s_n+u_{n-1}+s_{n-1}\\ &=2s_{n-1}+(s_{n-1}+2s_{n-2})+u_{n-1}+s_{n-1}\\ &=4s_{n-1}+2s_{n-2}+u_{n-1}.\tag{8} \end{align*}
But then by (5), and then substituting $4s_{n-3}=2u_{n-2}$ and $2s_{n-2}=u_{n-1}$ (by (6), and these are valid for $n\ge 6$), we find
 \begin{align*} 2s_{n-1}&=2s_{n-2}+4s_{n-3}\\ &=u_{n-1}+2u_{n-2}. \end{align*}
We substitute this into (8), finding that for $n\ge 6$,
 \begin{align*} a_{n}+a_{n-1}&=2s_{n-1}+2s_{n-2}+u_{n-1}+(u_{n-1}+2u_{n-2})\\ &=2(s_{n-1}+s_{n-2}+u_{n-1}+u_{n-2})\\ &=2(a_{n-1}+a_{n-2}).\tag{9} \end{align*}
Therefore, if $a_{n-1}+a_{n-2}=2^{n-1}$, then $a_{n}+a_{n-1}=2^n$. We already know that $a_4+a_3=2^4$. Also, we can compute that 
 \begin{align*} s_5&=s_4+2s_3=5+2\cdot 3=11\\ u_5&=2s_4=2\cdot 5=10. \end{align*}
Hence $a_5=s_5+u_5=21$. So as $a_4=s_4+u_4=11$, we find that $a_5+a_4=2^5$. Then we can use (9) for $n\ge 6$ to find by induction that $a_{n}+a_{n-1}=2^n$ for all $n\ge 4$.

 



\textbf{Solution 2}

We choose the tilings for the first and second pass around the circle simultaneously. As we tile both passes, our state is defined by whether or not we are in a continuation of a length 2 block on the first and second pass. We can't be in a continuation of a length 2 block on both passes, since that would mean we used the same move on the same block on both passes. Thus we have three states: (non-continuation, non-continuation), (non-continuation, continuation), and (continuation, non-continuation), where the elements of each tuple refer to the first and second pass respectively. Listing out the state transitions (while following the rule that we can't use the same move on both passes, and that a continuation must become a non-continuation), we get the transition matrix $T = \begin{bmatrix} 0 & 1 & 1 \\ 1 & 0 & 1 \\ 1 & 1 & 0 \end{bmatrix} = J-I$, where $J$ is the ones matrix, and $I$ is the identity. We have that $a_n = \left(T^n\right)_{11} + \left(T^n\right)_{23}$. The first term represents the case where we end the first pass at point $A$, and the second term represents the case where we jump over point $A$ when finishing the first pass. Thus it simply remains to compute $T^n$. We have

 \begin{align*}&\mathrel{\phantom{=}} T^n \\&= (J-I)^n \\&= \sum_{k=0}^n J^k (-1)^{n-k} {n \choose k} \\&= J \sum_{k=1}^n 3^{k-1} (-1)^{n-k} {n \choose k} + (-1)^n I \\&= \frac{1}{3}J\left(2^n-(-1)^n\right) + (-1)^n I\end{align*}


 Thus $a_n = \left(T^n\right)_{11} + \left(T^n\right)_{23} = \frac{2}{3}\left(2^n-(-1)^n\right) + (-1)^n$. Thus, $a_{n-1}+a_n = 2^n$, as desired.

 



\textbf{Solution 3}

We split this into two cases. Case 1: The marker lands on $A$ thrice. This implies that $A$ is touched after one full cycle. Place any nodes not touched in the first cycle in the set $K$. Lemma 1: Any unique configuration of $K$ such that [Work In Progress]~SigmaPiE

 


	\item (\textit{USAMO 2013}) Let $n$ be a positive integer.  There are $\frac{n(n+1)}{2}$ marks, each with a black side and a white side, arranged into an equilateral triangle, with the biggest row containing $n$ marks.  Initially, each mark has the black side up.  An \textit{operation} is to choose a line parallel to the sides of the triangle, and flipping all the marks on that line.  A configuration is called \textit{admissible} if it can be obtained from the initial configuration by performing a finite number of operations.  For each admissible configuration $C$, let $f(C)$ denote the smallest number of operations required to obtain $C$ from the initial configuration.  Find the maximum value of $f(C)$, where $C$ varies over all admissible configurations.


 



\textbf{Solution}

Solution by inconsistent 

 
This problem is headache-inducing... but also fun. Ouch.

 Let $n = 4k + b$ where $b \in {1, 2, 3, 4}$. Then I claim the answer is $6k+1, 6k+2, 6k+3, 6k+6$ in these four cases. First, note that the base cases $n = 1, 2, 3$ are trivial by checking to be $1, 2, 3$. Now note the case $n = 4$ needs at least $6$ to construct a black triangle with a single white mark at its center by checking. Call this size $4$ triangle the evil triangle.

 To show the lower bound, notice we can take the top $n$ rows of a $n+4$ equilateral triangle to be the limiting case for $n$ and put two evil triangles on the bottom left and bottom right, respectively: these force $4$ new rows to be occupied since each evil subcase requires at least two rows in every direction, which is a claim easy verified.

 To show the upper bound, we reparameterize the triangles in terms of their rows. Notably, we can start counting rows from each corner, and give each mark a barycentric coordinate $(a, b, c)$ where $a, b, c$ are integers denoting the $i$th row in a given direction the point lies in, and this satisfies $a+b+c = 2n+1$ by Viviani's theorem.

 Now, consider the configuration $C$ that minimizes $f(C)$. We place the rows in boxes: odd rows of each direction, and even rows of each direction, to make six boxes total. Call the numbers of odd rows $a, b, c$ and the even rows $d, e, f$ such that $(a, b, c)$ and $(d, e, f)$ correspond to the same directions respectively in that order. Notice no row is operated on twice since row operations are commutative. Assume for sake of contradiction our claim is wrong, then we can construct a solution for the case one higher than our claim.

 We now perform casework: for $b = 4$, it follows one of the quantities $a+b+d+e, b+c+e+f, c+a+f+d$ is greater than $n$. In this case, invert the rows chosen in those four boxes (i.e. remove the ones chosen and choose the ones unchosen). This fixes the configuration and decreases $f(C)$, giving contradiction.

 Now, consider the case $b = 1$. Like the above paragraph, we know $a+d, b+e, c+f$ are equal to $2k, 2k+1, 2k+1$ in some order since if any pair sums to more than $4k+1$ then we reach a contradiction. But then $2k+1 + 2k + 1 > 4k + 2$, giving contradiction.

 Now, for the two remaining cases, we show a lemma: the xor of all the odd rows in two directions and all the even row in a third direction, or the sum of all the even rows in all directions are both $0$. This result follows immediately from our earlier result via Vivani's theorem (it cannot be that all three claims are true, and it cannot be the case that two of the claims are false and one is true by parity).

 Thus in the $b = 2$ case it follows that $\min(a+b+f, b+c+d, c+a+e, d+e+f) \geq \frac{12k+6}{4}$ so one of the triples is at least $3k + 2$. However inverting the three involved boxes, that total to $6k+3$ rows and don't affect the configuration when all three are inverted, decreases the $f(C)$, giving contradiction.

 Finally, in the $b = 3$ it follows that $(a+b+f)+(b+c+d) + (c+a+e) + (d+e+f) =12k+8$. In particular, the odd/odd/even triples consist of $\frac{n+1}{2} + \frac{n+1}{2} + \frac{n-1}{2} = 6k+5$ rows each, while the even/even/even triple consists of $3 \cdot \frac{n-1}{2} = 6k+3$ rows. For each of the triples to avoid decreases $f(C)$ after being inverted, it needs to satisfy either $\leq 3k + 2$ or $\leq 3k+1$ depending on the two cases, so it follows that $LHS \leq 12k + 7$, giving contradiction. 

 Thus we have reached contradiction in all cases of counterexamples, proving our desired claim.

 


	\item (\textit{USAMO 2013}) Find all real numbers $x,y,z\geq 1$ satisfying 
 \[\min(\sqrt{x+xyz},\sqrt{y+xyz},\sqrt{z+xyz})=\sqrt{x-1}+\sqrt{y-1}+\sqrt{z-1}.\]


 



\textbf{Solution  (Cauchy or AM-GM)}

The key Lemma is:
 \[\sqrt{a-1}+\sqrt{b-1} \le \sqrt{ab}\] for all $a,b \ge 1$. Equality holds when $(a-1)(b-1)=1$.

 This is proven easily.
 \[\sqrt{a-1}+\sqrt{b-1} = \sqrt{a-1}\sqrt{1}+\sqrt{1}\sqrt{b-1} \le \sqrt{(a-1+1)(b-1+1)} = \sqrt{ab}\] by Cauchy.

 Equality then holds when $a-1 =\frac{1}{b-1} \implies (a-1)(b-1) = 1$.

 Now assume that $x = \min(x,y,z)$. Now note that, by the Lemma,

 
 \[\sqrt{x-1}+\sqrt{y-1}+\sqrt{z-1} \le \sqrt{x-1} + \sqrt{yz} \le \sqrt{x(yz+1)} = \sqrt{xyz+x}\]. So equality must hold in order for the condition in the problem statement to be met.So $(y-1)(z-1) = 1$ and $(x-1)(yz) = 1$. If we let $z = c$, then we can easily compute that $y = \frac{c}{c-1}, x = \frac{c^2+c-1}{c^2}$.Now it remains to check that $x \le y, z$.

 But by easy computations, $x = \frac{c^2+c-1}{c^2} \le c = z \Longleftrightarrow (c^2-1)(c-1) \ge 0$, which is obvious.Also $x = \frac{c^2+c-1}{c^2} \le \frac{c}{c-1} = y \Longleftrightarrow 2c \ge 1$, which is obvious, since $c \ge 1$.

 So all solutions are of the form $\boxed{\left(\frac{c^2+c-1}{c^2}, \frac{c}{c-1}, c\right)}$, and all permutations for $c > 1$.

 \textbf{Remark:} An alternative proof of the key Lemma is the following:By AM-GM, 
 \[(ab-a-b+1)+1 = (a-1)(b-1) + 1 \ge 2\sqrt{(a-1)(b-1)}\]
 \[ab\ge (a-1)+(b-1)+2\sqrt{(a-1)(b-1)}\]. Now taking the square root of both sides gives the desired. Equality holds when $(a-1)(b-1) = 1$.

 



\textbf{Solution 2}

WLOG,  assume that $x = \min(x,y,z)$. Let $a=\sqrt{x-1},$ $b=\sqrt{y-1}$ and $c=\sqrt{z-1}$. Then $x=a^2+1$, $y=b^2+1$ and $z=c^2+1$. The equation becomes
 \[(a^2+1)+(a^2+1)(b^2+1)(c^2+1)=(a+b+c)^2.\]Rearranging the terms, we have 
 \[(1+a^2)(bc-1)^2+[a(b+c)-1]^2=0.\]Therefore $bc=1$ and $a(b+c)=1.$ Express $a$ and $b$ in terms of $c$, we have $a=\frac{c}{c^2+1}$ and $b=\frac{1}{c}.$ Easy to check that $a$ is the smallest among $a$, $b$ and $c.$ Then $x=\frac{c^4+3c^2+1}{(c^2+1)^2}$, $y=\frac{c^2+1}{c^2}$ and $z=c^2+1.$Let $c^2=t$, we have the solutions for $(x,y,z)$ as follows:$(\frac{t^2+3t+1}{(t+1)^2}, \frac{t+1}{t}, t+1)$ and permutations for all $t>0.$ 

 --J.Z.

 


	\item (\textit{USAMO 2013}) Given positive integers $m$ and $n$, prove that there is a positive integer $c$ such that the numbers $cm$ and $cn$ have the same number of occurrences of each non-zero digit when written in base ten.


 



\textbf{Solution 1}

This solution is adopted from the \href{https://artofproblemsolving.comhttp://web.archive.org/web/20130606075948/http://amc.maa.org/usamo/2013/2013USAMO\_Day1\_Day2\_Final\_S.pdf}{official solution}. Both the problem and the solution were suggested by Richard Stong.

 Without Loss of Generality, suppose $m \geq n \geq 1$. By prime factorization of $n$, we can find a positive integer $c_1$ such that $c_1n=10^s n_1$ where $n_1$ is relatively prime to $10$. If a positive $k$ is larger than $s$, then $(10^k c_1 m - c_1 n)= 10^s t$, where $t=10^{k-s} c_1m-n_1>0$ is always relatively prime to $10$.

 Choose a $k$ large enough so that $t$ is larger than $c_1m$. We can find an integer $b\geq 1$ such that $10^b-1$ is divisible by $t$, and also larger than $10c_1m$. For example, let $b=\varphi(t)$ and use Euler's theorem. Now, let $c_2=(10^b-1)/t$, and $c=c_1c_2$. We claim that $c$ is the desired number.

 Indeed, since both $c_1m$ and $n_1$ are less than $t$, we see that the decimal expansion of both the fraction $(c_1m)/t = (cm)/(c_2t) = (cm)/(10^b-1)$ and $n_1/t=(c_2n_1)/(10^b-1)$ are repeated in $b$-digit. And we also see that $10^k (c_1m)/t = (t+c_1n)/t= 1+10^s (n_1/t)$, therefore the two repeated $b$-digit expansions are cyclic shift of one another.

 This proves that $cm$ and $c_2n_1$ have the same number of occurrences of non-zero digits. Furthermore, $cn = c_2c_1n=10^s c_2n_1$ also have the same number of occurrences of non-zero digits with $c_2n_1$.

 



\textbf{Solution 2}

This is a rephrasing of the above solution.

 It is enough to solve the problem when $m,n$ are replaced by $km,kn$ for any positive integer $k$. In particular, by taking $k=2^a5^b$ for appropriate values of $a,b$, we may assume $n=10^sn_1$ where $n_1$ is relatively prime to 10.

 Furthermore, adding or removing trailing zeros from $m$ and $n$ doesn't affect the claim, so we may further assume $\gcd(n,10)=1$ and that $m$ has a xillion trailing zeros (enough to make $m$ way bigger than $n$, and also so that $m$ has at least one trailing zero).

 For clarity of exposition, we will also multiply $m,n$ by a small number to make the units digit of $n$ be $1$ (though this is not necessary for the solution to work).

 The point is that, for any positive integer $X$, most nonzero digits appear the same number of times in $X$ and $X+999\dots999$ if there are enough $9$s; In particular, if the units digits of $X$ is $1$, then all nonzero digits appear the same number of times as long as there are at least as many $9$s as digits in $X$.

 So we will pick $c$ to satisfy:

 
                \begin{itemize}
                    \item $mc=nc+999\dots999$ where the number of $9$s is more than the number of digits of $nc$\item The units digit of $nc$ is $1$.
                \end{itemize}


                Because we made the units of $n$ to be $1$, the second condition is equivalent to making the units digit of $c$ to be $1$.

 The first condition is equivalent to $(m-n)c=999\dots999$. Because $m$ has at least one trailing 0, the units digit of $m-n$ is 9, so $\gcd(m-n,10)=1$ and there is some $c$ so that $(m-n)c=999\dots999$, and the units digit of $c$ must be $1$ which agrees with the other condition.

 Finally, as $m\gg n, (m-n)c\approx mc\gg nc$ so it is possible to make the number of $9$s more than the number of digits in $nc$.

 


 


	\item (\textit{USAMO 2013}) Let $ABC$ be a triangle.  Find all points $P$ on segment $BC$ satisfying the following property:  If $X$ and $Y$ are the intersections of line $PA$ with the common external tangent lines of the circumcircles of triangles $PAB$ and $PAC$, then 
 \[\left(\frac{PA}{XY}\right)^2+\frac{PB\cdot PC}{AB\cdot AC}=1.\]


 



\textbf{Solution}

Let circle $PAB$ (i.e. the circumcircle of $PAB$), $PAC$ be $\omega_1, \omega_2$ with radii $r_1$, $r_2$ and centers $O_1, O_2$, respectively, and $d$ be the distance between their centers.

 \textbf{Lemma.} $XY = \frac{r_1 + r_2}{d} \sqrt{d^2 - (r_1 - r_2)^2}.$

 Proof. Let the external tangent containing $X$ meet $\omega_1$ at $X_1$ and $\omega_2$ at $X_2$, and let the external tangent containing $Y$ meet $\omega_1$ at $Y_1$ and $\omega_2$ at $Y_2$. Then clearly $X_1 Y_1$ and $X_2 Y_2$ are parallel (for they are both perpendicular $O_1 O_2$), and so $X_1 Y_1 Y_2 X_2$ is a trapezoid.

 Now, $X_1 X^2 = XA \cdot XP = X_2 X^2$ by Power of a Point, and so $X$ is the midpoint of $X_1 X_2$. Similarly, $Y$ is the midpoint of $Y_1 Y_2$. Hence, $XY = \frac{1}{2} (X_1 Y_1 + X_2 Y_2).$ Let $X_1 Y_1$, $X_2 Y_2$ meet $O_1 O_2$ s at $Z_1, Z_2$, respectively. Then by similar triangles and the Pythagorean Theorem we deduce that $X_1 Z_1 = \frac{r_1 \sqrt{d^2 - (r_1 - r_2)^2}}{d}$ and $\frac{r_2 \sqrt{d^2 - (r_1 - r_2)^2}}{d}$. But it is clear that $Z_1$, $Z_2$ is the midpoint of $X_1 Y_1$, $X_2 Y_2$, respectively, so $XY = \frac{(r_1 + r_2)}{d} \sqrt{d^2 - (r_1 - r_2)^2},$ as desired.

 Lemma 2. Triangles $O_1 A O_2$ and $BAC$ are similar.

 Proof. $\angle{AO_1 O_2} = \frac{\angle{PO_1 A}}{2} = \angle{ABC}$ and similarly $\angle{AO_2 O_1} = \angle{ACB}$, so the triangles are similar by AA Similarity.

 Also, let $O_1 O_2$ intersect $AP$ at $Z$. Then obviously $Z$ is the midpoint of $AP$ and $AZ$ is an altitude of triangle $A O_1 O_2$.Thus, we can simplify our expression of $XY$:
 \[XY = \frac{AB + AC}{BC} \cdot \frac{AP}{2 h_a} \sqrt{BC^2 - (AB - AC)^2},\]where $h_a$ is the length of the altitude from $A$ in triangle $ABC$. Hence, substituting into our condition and using $AB = c, BC = a, CA = b$ gives
 \[\left( \frac{2a h_a}{(b+c) \sqrt{a^2 - (b-c)^2}} \right)^2 + \frac{PB \cdot PC}{bc} = 1.\]Using $2 a h_a = 4[ABC] = \sqrt{(a + b + c)(a + b - c)(a - b + c)(-a + b + c)}$ by Heron's Formula (where $[ABC]$ is the area of triangle $ABC$, our condition becomes
 \[\frac{(a + b + c)(-a + b + c)}{(b + c)^2} + \frac{PB \cdot PC}{bc} = 1,\]which by $(a + b + c)(-a + b + c) = (b + c)^2 - a^2$ becomes
 \[\frac{PB \cdot PC}{bc} = \frac{a^2 bc}{(b+c)^2}.\]Let $PB = x$; then $PC = a - x$. The quadratic in $x$ is
 \[x^2 - ax + \frac{a^2 bc}{(b+c)^2} = 0,\]which factors as
 \[\left(x - \frac{ab}{b+c}\right)\left(x - \frac{ac}{b+c}\right) = 0.\]Hence, $PB = \frac{ab}{b+c}$ or $\frac{ac}{b+c}$, and so the $P$ corresponding to these lengths are our answer.

 


\end{enumerate}